% GENERATED FILE. Edit canonical lessons, then run npm run generate:books.
% canonical-source-hash: fnv1a64:8cf3bf51939fa979
% canonical-lessons: KA-C263-huduga, KA-C263-hudugi, KA-C263-gandasu, KA-C263-hengasu, KA-C263-muduka

\chapter{Boy, Girl, Man, Woman, Old man}
\label{ch:ka-boy-girl-man-woman-old-man}
\hlchaptermodality{263}

\begin{chapteropening}
\textbf{By the end of this chapter:} \emph{I can say a boy, a girl, a man, a woman and an old man.}

Complete the last lesson of chapter 263: I can say a boy, a girl, a man, a woman and an old man.
\end{chapteropening}

\section[\texorpdfstring{\kn{ಹುಡುಗ} (huḍuga)}{ಹುಡುಗ (huḍuga)}]{\kn{ಹುಡುಗ} (huḍuga) — a boy}

\label{lesson:KA-C263-huduga}



\begin{quote}
Before the new one: say the Kannada for food offered to a god and then shared, then the Kannada for jewellery.
\end{quote}

\subsection*{You'll want to know: \kn{ಹುಡುಗ}}
\textbf{\kn{ಹುಡುಗ}} — \emph{huḍuga} — \textquotedblleft{}a boy\textquotedblright{}.

\begin{grammarlens}[title={What you've built}]
One more word for this chapter.
\end{grammarlens}

\subsection*{Your turn}
\begin{itemize}
\raggedright
  \item \emph{Say it:} \emph{huḍuga}
  \item \emph{Say it:} \emph{huḍuga}, once more
  \item \emph{Say it:} say \emph{huḍuga}
  \item \emph{Recall:} say the Kannada for food offered to a god and then shared, then the Kannada for jewellery, then say \emph{huḍuga} again
\end{itemize}

\begin{tcolorbox}[breakable,colback=teal!4,colframe=teal!35!black,title={Before you move on}]
What does \kn{ಹುಡುಗ} mean? (\textquotedblleft{}A boy\textquotedblright{}.) Say it once more.
\end{tcolorbox}
\section[\texorpdfstring{\kn{ಹುಡುಗಿ} (huḍugi)}{ಹುಡುಗಿ (huḍugi)}]{\kn{ಹುಡುಗಿ} (huḍugi) — a girl}

\label{lesson:KA-C263-hudugi}



\begin{quote}
Before the new one: say the Kannada for jewellery, then the Kannada for a boy.
\end{quote}

\subsection*{You'll want to know: \kn{ಹುಡುಗಿ}}
\textbf{\kn{ಹುಡುಗಿ}} — \emph{huḍugi} — \textquotedblleft{}a girl\textquotedblright{}.

\begin{grammarlens}[title={What you've built}]
One more word for this chapter.
\end{grammarlens}

\subsection*{Your turn}
\begin{itemize}
\raggedright
  \item \emph{Say it:} \emph{huḍugi}
  \item \emph{Say it:} \emph{huḍugi}, once more
  \item \emph{Say it:} say \emph{huḍugi}
  \item \emph{Recall:} say the Kannada for jewellery, then the Kannada for a boy, then say \emph{huḍugi} again
\end{itemize}

\begin{tcolorbox}[breakable,colback=teal!4,colframe=teal!35!black,title={Before you move on}]
What does \kn{ಹುಡುಗಿ} mean? (\textquotedblleft{}A girl\textquotedblright{}.) Say it once more.
\end{tcolorbox}
\section[\texorpdfstring{\kn{ಗಂಡಸು} (gaṇḍasu)}{ಗಂಡಸು (gaṇḍasu)}]{\kn{ಗಂಡಸು} (gaṇḍasu) — a man}

\label{lesson:KA-C263-gandasu}



\begin{quote}
Before the new one: say the Kannada for a boy, then the Kannada for a girl.
\end{quote}

\subsection*{You'll want to know: \kn{ಗಂಡಸು}}
\textbf{\kn{ಗಂಡಸು}} — \emph{gaṇḍasu} — \textquotedblleft{}a man\textquotedblright{}.

\begin{grammarlens}[title={What you've built}]
One more word for this chapter.
\end{grammarlens}

\subsection*{Your turn}
\begin{itemize}
\raggedright
  \item \emph{Say it:} \emph{gaṇḍasu}
  \item \emph{Say it:} \emph{gaṇḍasu}, once more
  \item \emph{Say it:} say \emph{gaṇḍasu}
  \item \emph{Recall:} say the Kannada for a boy, then the Kannada for a girl, then say \emph{gaṇḍasu} again
\end{itemize}

\begin{tcolorbox}[breakable,colback=teal!4,colframe=teal!35!black,title={Before you move on}]
What does \kn{ಗಂಡಸು} mean? (\textquotedblleft{}A man\textquotedblright{}.) Say it once more.
\end{tcolorbox}
\section[\texorpdfstring{\kn{ಹೆಂಗಸು} (heṅgasu)}{ಹೆಂಗಸು (heṅgasu)}]{\kn{ಹೆಂಗಸು} (heṅgasu) — a woman}

\label{lesson:KA-C263-hengasu}



\begin{quote}
Before the new one: say the Kannada for a girl, then the Kannada for a man.
\end{quote}

\subsection*{You'll want to know: \kn{ಹೆಂಗಸು}}
\textbf{\kn{ಹೆಂಗಸು}} — \emph{heṅgasu} — \textquotedblleft{}a woman\textquotedblright{}.

\begin{grammarlens}[title={What you've built}]
One more word for this chapter.
\end{grammarlens}

\subsection*{Your turn}
\begin{itemize}
\raggedright
  \item \emph{Say it:} \emph{heṅgasu}
  \item \emph{Say it:} \emph{heṅgasu}, once more
  \item \emph{Say it:} say \emph{heṅgasu}
  \item \emph{Recall:} say the Kannada for a girl, then the Kannada for a man, then say \emph{heṅgasu} again
\end{itemize}

\begin{tcolorbox}[breakable,colback=teal!4,colframe=teal!35!black,title={Before you move on}]
What does \kn{ಹೆಂಗಸು} mean? (\textquotedblleft{}A woman\textquotedblright{}.) Say it once more.
\end{tcolorbox}
\section[\texorpdfstring{\kn{ಮುದುಕ} (muduka)}{ಮುದುಕ (muduka)}]{\kn{ಮುದುಕ} (muduka) — an old man}

\label{lesson:KA-C263-muduka}



\begin{quote}
Before the new one: say the Kannada for a man, then the Kannada for a woman.
\end{quote}

\subsection*{You'll want to know: \kn{ಮುದುಕ}}
\textbf{\kn{ಮುದುಕ}} — \emph{muduka} — \textquotedblleft{}an old man\textquotedblright{}.

\begin{grammarlens}[title={What you've built}]
One more word for this chapter.
\end{grammarlens}

\subsection*{Your turn}
\begin{itemize}
\raggedright
  \item \emph{Say it:} \emph{muduka}
  \item \emph{Say it:} \emph{muduka}, once more
  \item \emph{Say it:} say \emph{muduka}
  \item \emph{Recall:} say the Kannada for a man, then the Kannada for a woman, then say \emph{muduka} again
\end{itemize}

\begin{tcolorbox}[breakable,colback=teal!4,colframe=teal!35!black,title={Before you move on}]
What does \kn{ಮುದುಕ} mean? (\textquotedblleft{}An old man\textquotedblright{}.) Say it once more.
\end{tcolorbox}
